\documentclass[12pt, a4paper, french]{article}

% --- Packages ---
\usepackage[utf8]{inputenc}
\usepackage[T1]{fontenc}
\usepackage{babel}
\usepackage{amsmath, amssymb, amsfonts}
\usepackage{graphicx}
\usepackage{geometry}
\usepackage{fancyhdr}
\usepackage{titlesec}
\usepackage{hyperref}
\usepackage{caption}
\usepackage{subcaption}
\usepackage{booktabs}
\usepackage{array}
\usepackage{float}
\usepackage{siunitx}
\usepackage{listings}
\usepackage{xcolor}
\usepackage{enumitem}
\usepackage{tocloft}
\usepackage[backend=bibtex, style=numeric]{biblatex}

% --- Configuration ---
\geometry{margin=2.5cm}
\addbibresource{references.bib}

\hypersetup{
    colorlinks=true,
    linkcolor=blue,
    citecolor=blue,
    urlcolor=blue,
    pdfauthor={Groupe XX},
    pdftitle={Problème III - Numérisation audio haute fidélité},
}

% --- En-têtes et pieds de page ---
\pagestyle{fancy}
\fancyhf{}
\fancyhead[L]{APP Signal - ISEP 2025-2026}
\fancyhead[R]{Problème III}
\fancyfoot[C]{\thepage}
\fancyfoot[L]{Groupe XX}
\renewcommand{\headrulewidth}{0.4pt}
\renewcommand{\footrulewidth}{0.4pt}

% --- Style des listings (code MATLAB) ---
\lstdefinestyle{matlab}{
    language=Matlab,
    basicstyle=\small\ttfamily,
    keywordstyle=\color{blue}\bfseries,
    stringstyle=\color{red},
    commentstyle=\color{green!60!black}\itshape,
    numbers=left,
    numberstyle=\tiny\color{gray},
    numbersep=5pt,
    frame=single,
    breaklines=true,
    tabsize=4,
    showstringspaces=false,
}

% --- Commandes personnalisées ---
\newcommand{\Fe}{F_e}
\newcommand{\fmax}{f_{\text{max}}}
\newcommand{\fmin}{f_{\text{min}}}
\newcommand{\SNR}{\text{SNR}}
\newcommand{\dB}{\text{dB}}
\newcommand{\dBm}{\text{dBm}}

% =====================================================================
\begin{document}

% --- PAGE DE GARDE ---
\begin{titlepage}
    \centering
    \vspace*{2cm}
    
    {\Large \textbf{ISEP -- Institut Supérieur d'Électronique de Paris}}\\[0.5cm]
    {\large APP Composante Signal -- 2025-2026}\\[2cm]
    
    \rule{\textwidth}{1pt}\\[0.5cm]
    {\LARGE \textbf{Problème III}}\\[0.3cm]
    {\Large \textbf{Numérisation d'un signal audio haute fidélité}}\\[0.5cm]
    \rule{\textwidth}{1pt}\\[2cm]
    
    {\large
    \begin{tabular}{ll}
        \textbf{Groupe}   & : XX \\[0.3cm]
        \textbf{Membres}  & : Nom Prénom 1 \\
                          & \phantom{:} Nom Prénom 2 \\
                          & \phantom{:} Nom Prénom 3 \\[0.3cm]
        \textbf{Tuteur}   & : Nom du tuteur \\[0.3cm]
        \textbf{Date}     & : \today \\
    \end{tabular}
    }
    
    \vfill
    {\small Livrable II -- Rapport sur la résolution du Problème III}
\end{titlepage}

% --- TABLE DES MATIÈRES ---
\tableofcontents
\newpage

% =====================================================================
% INTRODUCTION
% =====================================================================
\section{Introduction}

L'objectif de ce rapport est de concevoir un système de numérisation d'un signal audio en haute fidélité. Le cahier des charges est le suivant :

\begin{itemize}
    \item Bande de fréquences : de \SI{20}{\hertz} à \SI{20000}{\hertz}
    \item Rapport signal à bruit (SNR) moyen : \SI{90}{\dB}
    \item Microphone idéal (bande infinie, aucune distorsion, aucun bruit)
    \item Dynamique du microphone : $[-\SI{500}{\milli\volt}, +\SI{500}{\milli\volt}]$
    \item Puissance moyenne en sortie du micro : \SI{30}{\milli\watt}
\end{itemize}

Nous devons déterminer le schéma fonctionnel complet de la chaîne de numérisation, calculer tous les paramètres nécessaires (fréquence d'échantillonnage, nombre de bits de quantification, conception du filtre anti-repliement), puis évaluer la capacité de stockage requise et la compatibilité avec le standard CD audio.

Ce rapport présente d'abord le schéma fonctionnel de la chaîne de numérisation (section~\ref{sec:schema}), puis le calcul détaillé de chaque paramètre (section~\ref{sec:parametres}), suivi de la conception du filtre anti-repliement (section~\ref{sec:filtre}). Enfin, nous analysons la compatibilité avec le standard CD et le respect du timbre des instruments d'un orchestre symphonique (section~\ref{sec:compatibilite}).

% =====================================================================
% SCHÉMA FONCTIONNEL
% =====================================================================
\section{Schéma fonctionnel de la chaîne de numérisation}
\label{sec:schema}

La numérisation d'un signal analogique repose sur trois étapes fondamentales, précédées d'un filtrage anti-repliement. Le schéma fonctionnel est présenté en Figure~\ref{fig:schema}.

\begin{figure}[H]
    \centering
    % \includegraphics[width=\textwidth]{figures/schema_numerisation.png}
    \fbox{\parbox{0.9\textwidth}{\centering\vspace{1cm}[Insérer ici le schéma fonctionnel]\\\vspace{0.5cm}
    Micro $\rightarrow$ Filtre anti-repliement $\rightarrow$ Échantillonneur $\rightarrow$ Quantificateur $\rightarrow$ Codeur\vspace{1cm}}}
    \caption{Schéma fonctionnel de la chaîne de numérisation}
    \label{fig:schema}
\end{figure}

Les blocs sont les suivants :
\begin{enumerate}
    \item \textbf{Microphone} : capteur idéal, convertit le signal acoustique en signal électrique $x_c(t)$.
    \item \textbf{Filtre anti-repliement} : filtre passe-bas de fréquence de coupure $f_c$, élimine les composantes spectrales au-delà de $\Fe/2$ pour garantir le respect du théorème de Shannon.
    \item \textbf{Échantillonneur} : prélève les valeurs du signal filtré à la fréquence $\Fe$, produisant $x_e(n T_e)$.
    \item \textbf{Quantificateur (CAN)} : code les amplitudes sur $b$ bits avec un pas de quantification $q$.
    \item \textbf{Codeur} : met en forme le signal numérique pour le stockage ou la transmission.
\end{enumerate}

% =====================================================================
% CALCUL DES PARAMÈTRES
% =====================================================================
\section{Calcul des paramètres de numérisation}
\label{sec:parametres}

\subsection{Fréquence d'échantillonnage}

D'après le \textbf{théorème de Shannon}, pour échantillonner un signal sans perte d'information, il faut :
\begin{equation}
    \Fe \geq 2 \cdot \fmax = 2 \times \SI{20000}{\hertz} = \SI{40000}{\hertz}
    \label{eq:shannon}
\end{equation}

En pratique, une marge est nécessaire pour permettre la réalisation du filtre anti-repliement (bande de transition non nulle). Le standard CD Audio utilise $\Fe = \SI{44100}{\hertz}$, ce qui offre :
\begin{itemize}
    \item $\Fe/2 = \SI{22050}{\hertz} > \fmax = \SI{20000}{\hertz}$ : Shannon respecté.
    \item Bande de transition : $\Delta f = \Fe/2 - \fmax = \SI{2050}{\hertz}$, suffisante pour un filtre d'ordre raisonnable.
\end{itemize}

\begin{equation}
    \boxed{\Fe = \SI{44100}{\hertz}}
\end{equation}

\subsection{Nombre de bits de quantification}

Le pas de quantification est $q = A / 2^b$, avec $A = 2 V_{\text{max}} = \SI{1}{\volt}$ la dynamique totale.

La puissance du bruit de quantification est :
\begin{equation}
    P_e = \frac{q^2}{12} = \frac{A^2}{12 \cdot 2^{2b}}
    \label{eq:pe}
\end{equation}

Le rapport signal à bruit de quantification s'écrit :
\begin{equation}
    \SNR = 10 \log_{10}\left(\frac{P_{\text{moy}}}{P_e}\right) = 10 \log_{10}\left(\frac{12 \cdot P_{\text{moy}}}{A^2}\right) + 20 \cdot b \cdot \log_{10}(2)
    \label{eq:snr}
\end{equation}

En substituant $P_{\text{moy}} = \SI{30}{\milli\watt}$ et $A = \SI{1}{\volt}$ :
\begin{equation}
    \SNR = 10 \log_{10}\left(\frac{12 \times 0.030}{1^2}\right) + 6.02 \cdot b = -4.44 + 6.02 \cdot b \quad (\dB)
\end{equation}

Pour atteindre $\SNR \geq \SI{90}{\dB}$ :
\begin{equation}
    b \geq \frac{90 - (-4.44)}{6.02} = \frac{94.44}{6.02} \approx 15.69 \quad \Rightarrow \quad b = 16 \text{ bits}
\end{equation}

\textbf{Vérification} : avec $b = 16$ bits :
\begin{align}
    q &= \frac{1}{2^{16}} = \SI{15.26}{\micro\volt} \\
    P_e &= \frac{q^2}{12} = \SI{1.94e-11}{\watt} \\
    \SNR &= 10 \log_{10}\left(\frac{0.030}{1.94 \times 10^{-11}}\right) = \SI{91.90}{\dB} > \SI{90}{\dB} \quad \checkmark
\end{align}

\begin{equation}
    \boxed{b = 16 \text{ bits}, \quad \SNR = \SI{91.90}{\dB}}
\end{equation}

\subsection{Débit binaire}

Le débit binaire pour un enregistrement stéréo est :
\begin{equation}
    D = \Fe \times b \times n_{\text{canaux}} = 44100 \times 16 \times 2 = \SI{1411200}{bits/s} \approx \SI{1411.2}{kbits/s}
    \label{eq:debit}
\end{equation}

\begin{equation}
    \boxed{D = \SI{1411.2}{kbits/s} \approx \SI{1.41}{Mbits/s}}
\end{equation}

\subsection{Capacité de stockage}

Pour une heure de concert stéréo ($T = \SI{3600}{\second}$) :
\begin{equation}
    V = D \times T = 1\,411\,200 \times 3600 = 5\,080\,320\,000 \text{ bits} = 635\,040\,000 \text{ octets} \approx \SI{605.6}{\mega\byte}
\end{equation}

\begin{equation}
    \boxed{V \approx \SI{605.6}{\mega\byte} \approx \SI{0.59}{\giga\byte}}
\end{equation}

% =====================================================================
% FILTRE ANTI-REPLIEMENT
% =====================================================================
\section{Conception du filtre anti-repliement}
\label{sec:filtre}

\subsection{Gabarit du filtre}

Le filtre anti-repliement est un filtre passe-bas dont les paramètres sont déterminés par le contexte :

\begin{table}[H]
    \centering
    \caption{Paramètres du gabarit du filtre anti-repliement}
    \label{tab:gabarit}
    \begin{tabular}{lll}
        \toprule
        \textbf{Paramètre} & \textbf{Symbole} & \textbf{Valeur} \\
        \midrule
        Fréquence de coupure       & $f_c$        & \SI{20000}{\hertz} \\
        Début bande atténuée       & $f_a$        & \SI{22050}{\hertz} \\
        Bande de transition        & $\Delta f$   & \SI{2050}{\hertz} \\
        Ondulation bande passante  & $A_{\max}$   & \SI{0.1}{\dB} \\
        Atténuation bande atténuée & $A_{\min}$   & \SI{90}{\dB} \\
        \bottomrule
    \end{tabular}
\end{table}

\subsection{Justification des paramètres}

\begin{itemize}
    \item $f_c = \SI{20000}{\hertz}$ : fréquence maximale audible, limite de la bande utile.
    \item $f_a = \Fe/2 = \SI{22050}{\hertz}$ : au-delà de cette fréquence, les composantes seraient repliées.
    \item $A_{\min} = \SI{90}{\dB}$ : l'atténuation doit être au moins égale au SNR cible pour que le bruit de repliement soit négligeable devant le bruit de quantification.
\end{itemize}

\subsection{Résultats de synthèse}

Le filtre a été synthétisé sous MATLAB/Python par la méthode de Parks-McClellan (equiripple). Les résultats sont présentés en Figure~\ref{fig:filtre}.

\begin{figure}[H]
    \centering
    % \includegraphics[width=0.9\textwidth]{figures/filtre_anti_repliement.png}
    \fbox{\parbox{0.8\textwidth}{\centering\vspace{2cm}[Insérer ici les courbes du filtre]\vspace{2cm}}}
    \caption{Réponse en fréquence du filtre anti-repliement}
    \label{fig:filtre}
\end{figure}

% =====================================================================
% COMPATIBILITÉ
% =====================================================================
\section{Compatibilité CD et respect du timbre}
\label{sec:compatibilite}

\subsection{Compatibilité avec le standard CD Audio}

Le standard CD Audio (Red Book, 1982) spécifie :
\begin{itemize}
    \item $\Fe = \SI{44100}{\hertz}$, $b = 16$ bits, stéréo
    \item Capacité : \SI{700}{\mega\byte}, durée maximale $\approx \SI{80}{min}$
\end{itemize}

Nos paramètres ($\Fe = \SI{44100}{\hertz}$, $b = 16$ bits, stéréo) sont \textbf{identiques} au standard CD. Le volume de stockage pour 1~heure (\SI{605.6}{\mega\byte}) est inférieur à la capacité d'un CD (\SI{700}{\mega\byte}).

\begin{equation}
    \boxed{\text{Compatible CD} : \SI{605.6}{\mega\byte} < \SI{700}{\mega\byte} \quad \checkmark}
\end{equation}

\subsection{Respect du timbre des instruments}

Le timbre d'un instrument est déterminé par le contenu en harmoniques de chaque note jouée. Un orchestre symphonique comprend des instruments dont les harmoniques s'étendent jusqu'à environ \SI{16000}{\hertz} (cymbales) voire \SI{20000}{\hertz}.

Avec $\Fe/2 = \SI{22050}{\hertz}$, toutes les harmoniques audibles sont capturées. Le SNR de \SI{91.90}{\dB} garantit une dynamique suffisante pour restituer les nuances d'un orchestre (dynamique typique : \SI{60}{\dB} à \SI{80}{\dB}).

\begin{equation}
    \boxed{\text{Timbre respecté : toutes les harmoniques audibles sont capturées} \quad \checkmark}
\end{equation}

% =====================================================================
% CONCLUSION
% =====================================================================
\section{Conclusion}

La conception du système de numérisation audio haute fidélité aboutit aux paramètres suivants :

\begin{table}[H]
    \centering
    \caption{Récapitulatif des paramètres de numérisation}
    \label{tab:recap}
    \begin{tabular}{ll}
        \toprule
        \textbf{Paramètre} & \textbf{Valeur} \\
        \midrule
        Fréquence d'échantillonnage & \SI{44100}{\hertz} \\
        Nombre de bits               & 16 bits \\
        Pas de quantification        & \SI{15.26}{\micro\volt} \\
        SNR de quantification        & \SI{91.90}{\dB} \\
        Débit stéréo                 & \SI{1411.2}{kbits/s} \\
        Stockage 1h stéréo           & \SI{605.6}{\mega\byte} \\
        \bottomrule
    \end{tabular}
\end{table}

Ces paramètres sont identiques au standard CD Audio, confirmant que ce format a été conçu précisément pour l'enregistrement audio haute fidélité. Le stockage d'une heure de concert stéréo (\SI{605.6}{\mega\byte}) est compatible avec la capacité d'un CD (\SI{700}{\mega\byte}). Le timbre des instruments d'un orchestre symphonique est intégralement respecté, toutes les harmoniques audibles étant capturées dans la bande $[0, \SI{22050}{\hertz}]$.

Des améliorations pourraient être envisagées avec des formats de plus haute résolution (par exemple DVD-Audio : \SI{96}{\kilo\hertz}/24 bits), offrant un SNR supérieur et une bande passante étendue, au prix d'un débit et d'un stockage plus importants.

% =====================================================================
% RÉFÉRENCES
% =====================================================================
\newpage
\printbibliography[title={Références}]

% =====================================================================
% LISTE DES FIGURES ET TABLEAUX
% =====================================================================
\newpage
\listoffigures
\listoftables

% =====================================================================
% ANNEXES
% =====================================================================
\newpage
\appendix
\section{Annexe A : Démonstration du SNR de quantification}

Le bruit de quantification $e(nT_e) = x_e(nT_e) - x_q(nT_e)$ est modélisé comme une variable aléatoire uniformément distribuée sur $[-q/2, +q/2]$, où $q$ est le pas de quantification.

La puissance (variance) de ce bruit est :
\begin{equation}
    P_e = \mathbb{E}[e^2] = \int_{-q/2}^{q/2} e^2 \cdot \frac{1}{q} \, de = \frac{1}{q} \left[\frac{e^3}{3}\right]_{-q/2}^{q/2} = \frac{q^2}{12}
\end{equation}

Avec $q = A/2^b$ :
\begin{equation}
    P_e = \frac{A^2}{12 \cdot 2^{2b}}
\end{equation}

Le SNR de quantification en dB est :
\begin{equation}
    \SNR = 10\log_{10}\left(\frac{P_s}{P_e}\right) = 10\log_{10}(P_s) - 10\log_{10}\left(\frac{A^2}{12 \cdot 2^{2b}}\right)
\end{equation}

Pour un signal sinusoïdal pleine échelle ($P_s = A^2/8$) :
\begin{equation}
    \SNR = 10\log_{10}\left(\frac{A^2/8}{A^2/(12 \cdot 2^{2b})}\right) = 10\log_{10}\left(\frac{3}{2} \cdot 2^{2b}\right) = 6.02 \cdot b + 1.76 \; \dB
\end{equation}

\section{Annexe B : Codes MATLAB}

\textit{Les codes MATLAB sont fournis dans le dossier \texttt{src/} du fichier zip joint.}

% Mode d'emploi
\section{Annexe C : Mode d'emploi}

\subsection{Exécution des scripts MATLAB}
\begin{enumerate}
    \item Ouvrir MATLAB
    \item Se placer dans le dossier \texttt{src/}
    \item Exécuter \texttt{main\_probleme3.m} pour la résolution complète
    \item Les scripts individuels peuvent être exécutés séparément
\end{enumerate}

\subsection{Exécution des scripts Python}
\begin{enumerate}
    \item Installer les dépendances : \texttt{pip install numpy scipy matplotlib}
    \item Se placer dans le dossier \texttt{python/}
    \item Exécuter \texttt{python main\_probleme3.py}
\end{enumerate}

\end{document}
